\subsection{Сила трения}
%%%%%%%%%%%%%%%%%%%%%%%%%%%%%%%%%%%%%%%%%%%%%%%%%%%%
%%%%%%%%%%%%%%%%%%%%%%%%%%%%%%%%%%%%%%%%%%%%%%%%%%%%
\z{На столе лежат в стопке три разных бруска. Чтобы сдвинуть верхний, надо приложить силу $7 \,\text{Н}$, а чтобы вытянуть средний, придерживая (но не приподнимая) верхний, надо приложить силу $24 \,\text{Н}$. Какую силу надо приложить, чтобы сдвинуть два верхних бруска вместе?
}
%
{$17 \,\text{Н}$}

%%%%%%%%%%%%%%%%%%%%%%%%%%%%%%%%%%%%%%%%%%%%%%%%%%%%
\z[{\begin{wrapfigure}[3]{r}{0.3\textwidth}%
		\vspace{-15pt}%
		\begin{flushright}%
 			\includegraphics[width=0.3\textwidth]{3.5.2}%
		\end{flushright}%
	\end{wrapfigure}%
	}]%
{Два груза массой $3m$ и $m$, связанные невесомой нерастяжимой нитью, находятся на горизонтальном столе. Коэффициент трения грузов о стол равен $\mu$.
	\lpu Какую минимальную горизонтальную силу $F_{\min}$ нужно приложить к грузу массой m, чтобы система пришла в движение? \lpuend 
	\lpu С каким ускорением будут двигаться грузы в случае, если приложить к грузу $m$ силу $F>F_{\min}$?\lpuend 
	\lpu Чему в этом случае будет равна сила натяжения нити?  \lpuend 
}
%
{\puan $F_{\min}=4\mu mg$;\quad\puan $a=\frac{F-4\mu mg}{4m}$;\quad\puan $T=\frac{3F}{4}$}


%%%%%%%%%%%%%%%%%%%%%%%%%%%%%%%%%%%%%%%%%%%%%%%%%%%%
\z[{\begin{wrapfigure}[4]{r}{0.3\textwidth}%
		\vspace{-10pt}%
		\begin{flushright}%
 			\includegraphics[width=0.3\textwidth]{3.5.3}%
		\end{flushright}%
	\end{wrapfigure}%
	}]%
{Два ящика покоятся на горизонтальной поверхности, как показано на рисунке. Чтобы преодолеть трение и сдвинуть с места левый ящик, к нему необходимо приложить горизонтальную силу больше $F_1 = 24 \,\text{Н}$. Чтобы сдвинуть правый — $F_2 = 18 \,\text{Н}$. Ящики соединили нитью, переброшенной через блоки. Какую минимальную силу надо приложить к концу нити, чтобы расстояние между ящиками начало уменьшаться?
}
%
{$8 \,\text{Н}$}

%%%%%%%%%%%%%%%%%%%%%%%%%%%%%%%%%%%%%%%%%%%%%%%%%%%%
\z{На тело массы $m$, вначале покоившееся на горизонтальной плоскости, в течение времени $\tau$ действует горизонтальная сила $F$. Какое расстояние $L$ пройдет тело за время движения? Коэффициент трения тела о плоскость $\mu$.
}
%
{$L = \frac{(F - \mu m g) F \tau^2}{2 \mu m^2 g}$}

%%%%%%%%%%%%%%%%%%%%%%%%%%%%%%%%%%%%%%%%%%%%%%%%%%%%
\z[{\begin{wrapfigure}[3]{r}{0.3\textwidth}%
		\vspace{-10pt}%
		\begin{flushright}%
 			\includegraphics[width=0.3\textwidth]{3.5.5}%
		\end{flushright}%
	\end{wrapfigure}%
	}]%
{На горизонтальном шероховатом столе помещены грузы $M$ (внизу) и $m$, связанные нитью, переброшенной через неподвижный блок. Коэффициенты трения грузов друг о друга и нижнего груза о поверхность стола одинаковы и равны $\mu$. С какой горизонтальной силой нужно потянуть нижний груз, чтобы тела пришли в движение? Нити горизонтальны.
}
%
{$F = (3m + M) \mu g$}

%%%%%%%%%%%%%%%%%%%%%%%%%%%%%%%%%%%%%%%%%%%%%%%%%%%%
\z[{\begin{wrapfigure}[3]{r}{0.3\textwidth}%
		\vspace{-15pt}%
		\begin{flushright}%
 			\includegraphics[width=0.3\textwidth]{3.5.6}%
		\end{flushright}%
	\end{wrapfigure}%
	}]%
{На горизонтальной гладкой поверхности стола покоится доска массой $M$ (см. рисунок). На доску со скоростью $v$ въезжает шайба массы $m$. Какой должна быть длина доски, чтобы шайба не соскользнула с неё? Коэффициент трения скольжения между шайбой и доской равен $\mu$, размер шайбы мал по сравнению с длиной доски.
}
%
{$L = \frac{M v^2}{2 \mu (M + m) g}$}

%%%%%%%%%%%%%%%%%%%%%%%%%%%%%%%%%%%%%%%%%%%%%%%%%%%%
\z[{\begin{wrapfigure}[3]{r}{0.3\textwidth}%
		\vspace{-15pt}%
		\begin{flushright}%
 			\includegraphics[width=0.3\textwidth]{3.5.7}%
		\end{flushright}%
	\end{wrapfigure}%
	}]%
{На гладком горизонтальном столе расположена система грузов, изображённая на рисунке. Правый нижний груз тянут вдоль стола с силой $F$, как указано на рисунке. Коэффициент трения между грузами массы $m_1$ и $m_2$ равен $\mu$. Найдите ускорение всех грузов системы.
}
%
{При $F<\frac{2\mu m_1g(m_1+m_2)}{m_2+2m_1}  \equiv F_0$: $a_{1\text{лев}} = a_{1\text{прав}} = a_{2\text{прав}} = \frac{F}{2(m_1 + m_2)}$
\par
При $F>F_0$: $a_{2\text{прав}} = \frac{F - \mu m_1 g}{m_2}, \quad a_{1\text{лев}} = a_{1\text{прав}} = a_{2\text{лев}} = \frac{\mu m_1 g}{m_2 + 2 m_1}$
}